\section{The large-concentration limit of the direct jump}\label{app:hdp}

We prove Proposition~\ref{prop:directjump}(b), with $B_1,B_2,B_3$ independent
$\mathrm{Beta}(a,a)$ as in part (a) and $a=\omega/3\ge1$.

\emph{Idea of the proof.} We pass to the logits $\ell_j$ of the $B_j$. In
these coordinates the three conditions that define $D$ say that three linear
forms $M_1,M_2,M_3$ of $\ell$, whose sum is $0$, are bounded above by
quantities of second order in $\ell$. For large $a$ the law of $\ell$
concentrates near $0$ on the scale $a^{-1/2}$. There the bounds are small, so
on $D$ all three forms are nearly $0$. The points where they vanish form the
line of the cyclic kernels, so $D$ is a thin tube around this line. Let $t$ be
the coordinate along the line. At the point $t$ of the line each of the three
bounds is about $\frac12t^2$, so the section of $D$ at $t$ is close to the
triangle $M_1,M_2,M_3\le\frac12t^2$, which has area $\frac98t^4$. So we expect
$\Prob_\omega(D)\approx a^{3/2}\int t^4e^{-cat^2}dt\approx1/a$, and a Gaussian
integral gives the constant.

Before the proof we set up three things, namely the logits and the form of $D$
in them, two bounds on the function that appears, and coordinates along and
across the cyclic kernels.

\emph{Logits.} Put $\ell_j=\log(B_j/(1-B_j))$, $g(y)=2\log\cosh(y/2)$ and
$g_j=g(\ell_j)$. Since $B_j=e^{\ell_j/2}/(2\cosh(\ell_j/2))$ and
$1-B_j=e^{-\ell_j/2}/(2\cosh(\ell_j/2))$, we have
\[
 \log(2B_j)=\tfrac12(\ell_j-g_j),\qquad
 \log\bigl(2(1-B_j)\bigr)=-\tfrac12(\ell_j+g_j),\qquad
 B_j(1-B_j)=\tfrac14e^{-g_j}.
\]
The function $g$ is even and convex with $g(0)=g'(0)=0$, so $g\ge0$. The
density of $B_j$ is proportional to $(B_j(1-B_j))^{a-1}$, and
$dB_j/d\ell_j=B_j(1-B_j)$. So $\ell_j$ has the density $e^{-ag(\ell_j)}/Z$ with
$Z=4^a\Gamma(a)^2/\Gamma(2a)=2\sqrt\pi\,\Gamma(a)/\Gamma(a+\frac12)$, where the
second form is the duplication formula for $\Gamma$.

By the proof of Proposition~\ref{prop:directjump}, $D$ is the event that
$B_1B_2$, $(1-B_1)B_3$ and $(1-B_2)(1-B_3)$ are all at most $\frac14$. We
multiply each product by $4$ and take logarithms with the display above. For
example the first condition becomes
$\frac12(\ell_1-g_1)+\frac12(\ell_2-g_2)\le0$. So $D$ is the set where
\begin{equation}\label{eq:hdp-constraints}
 M_1\le g_1+g_2,\qquad M_2\le g_1+g_3,\qquad M_3\le g_2+g_3,
\end{equation}
with $M_1=\ell_1+\ell_2$, $M_2=\ell_3-\ell_1$ and
$M_3=-\ell_2-\ell_3=-M_1-M_2$. Let $\psi=g_1+g_2+g_3$ (in this appendix $\psi$
is not the Bessel function of Section~\ref{sec:planar}). Then the joint density
of $\ell=(\ell_1,\ell_2,\ell_3)$ is $Z^{-3}e^{-a\psi(\ell)}$, and
\begin{equation}\label{eq:hdp-integral}
 \Prob_\omega(D)=Z^{-3}\int_De^{-a\psi(\ell)}d\ell.
\end{equation}

\emph{Two bounds on $g$.} For real $y$,
\begin{align}
 \frac{y^2}4-\frac{y^4}{96}\le g(y)&\le\frac{y^2}4,\tag{g1}\label{eq:hdp-g1}\\
 g(y)&\ge\frac1{10}\min(y^2,|y|).\tag{g2}\label{eq:hdp-g2}
\end{align}
The first says that $g(y)$ is close to $y^2/4$ near $0$, and we use it near
the cyclic kernels. The second is a rough lower bound that we use far from
$0$. For \eqref{eq:hdp-g1} put $x=y/2$, so that $g(y)=2\log\cosh x$. The upper
bound is $\cosh x\le e^{x^2/2}$, which holds term by term in the two power
series. The lower bound says that
$\chi(x)=\log\cosh x-\frac{x^2}2+\frac{x^4}{12}$ is nonnegative, and this holds
as $\chi(0)=\chi'(0)=0$ and $\chi''=x^2-\tanh^2x\ge0$. For \eqref{eq:hdp-g2},
$g''=\frac12\mathrm{sech}^2(y/2)>0.2$ on $|y|\le2$, so $g(y)\ge0.1y^2$ there.
For $|y|\ge2$ convexity and $g(0)=0$ give $g(y)\ge\frac12|y|g(2)>0.43|y|$.

\emph{Coordinates along and across the cyclic kernels.} A cyclic kernel has
$B_1=1-B_2=B_3$, so its logits are $\ell=tv$ with $v=(1,-1,1)$ and $t\in\R$,
and $M_1=M_2=M_3=0$ there. For general $\ell$ let $t=\langle\ell,v\rangle/3$.
We use $(t,M_1,M_2)$ as coordinates. The map $\ell\mapsto(t,M_1,M_2)$ is
linear with determinant $1$, and its inverse is $\ell=tv+w$ with
$w=\frac13(M_1-M_2,2M_1+M_2,M_1+2M_2)\perp v$. So $t$ measures the position
along the line and $w$ the displacement from it. Since $|v|^2=3$,
$|\ell|^2=3t^2+|w|^2$, and a direct computation gives
$|w|^2=\frac13(M_1^2+M_2^2+M_3^2)=\frac23(M_1^2+M_1M_2+M_2^2)$. Write $|M|$
and $|M|_\infty$ for the Euclidean norm and the maximum norm of $(M_1,M_2)$.
As $\lvert M_1M_2\rvert\le\frac12|M|^2$,
\begin{equation}\label{eq:hdp-w}
 \tfrac13|M|^2\le|w|^2\le|M|^2\le2|M|_\infty^2.
\end{equation}

\begin{proof}[Proof of Proposition~\ref{prop:directjump}(b)]
We compute the integral in \eqref{eq:hdp-integral} in four steps. Step 1 shows
that on $D$ the forms $M_i$ are of second order in $\ell$. Step 2 removes the
part of $D$ away from $0$. Step 3 finds the section of $D$ at $t$, and Step 4
integrates over $t$.

\emph{Step 1: an a priori bound.} On $D$, each right side
in \eqref{eq:hdp-constraints} is at most $\psi$, so $M_1\le\psi$, and
$M_1=-M_2-M_3\ge-2\psi$. The same holds for $M_2$. So by \eqref{eq:hdp-g1},
which gives $\psi\le\frac14|\ell|^2$,
\begin{equation}\label{eq:hdp-apriori}
 |M|_\infty\le2\psi(\ell)\le\tfrac12|\ell|^2=\tfrac12(3t^2+|w|^2).
\end{equation}

\emph{Step 2: localization.} Put $\varepsilon=1/12$ and let $U$ be the set
where $|t|\le\varepsilon$ and $|M|_\infty\le\varepsilon$. We show that the
integral over $D\setminus U$ is exponentially small in $a$. Off $U$ we have
$|t|>\varepsilon$ or $|M|>\varepsilon$, so by \eqref{eq:hdp-w}
$|\ell|\ge\max(\sqrt3|t|,|M|/\sqrt3)\ge\varepsilon/\sqrt3$. Also some
coordinate has $|\ell_j|\ge|\ell|/\sqrt3$. So \eqref{eq:hdp-g2} gives
\[
 \psi\ge\gamma:=\tfrac1{10}\min\bigl(|\ell|^2/3,\,|\ell|/\sqrt3\bigr)
 \ge\gamma_0:=\varepsilon^2/90 .
\]
Hence $a\psi\ge(a-1)\gamma_0+\gamma$, that is,
$e^{-a\psi}\le e^{-(a-1)\gamma_0}e^{-\gamma}$ off $U$. The function
$e^{-\gamma(\ell)}$ is integrable over $\R^3$, so
$\int_{D\setminus U}e^{-a\psi}=O(e^{-\gamma_0a})$.

\emph{Step 3: the section of $D$ at $t$.} We show that the section of
$D\cap U$ at $t$ lies between two triangles of area close to $\frac98t^4$.
First, on $D\cap U$ the bounds \eqref{eq:hdp-apriori} and \eqref{eq:hdp-w} give
\[
 |M|_\infty\le\tfrac32t^2+\tfrac12|w|^2\le\tfrac32t^2+|M|_\infty^2
 \le\tfrac32t^2+\varepsilon|M|_\infty,
\]
so $|M|_\infty\le\frac32t^2/(1-\varepsilon)\le2t^2$. So $D\cap U$ lies in the
set $E=\{|t|\le\varepsilon,\ |M|_\infty\le2t^2\}$.

Next we estimate $g_j$ and $\psi$ on $E$. There $|w|^2\le8t^4$ by
\eqref{eq:hdp-w}, and $\ell_j=\pm t+w_j$ with $|w_j|\le|w|$. So
$|\ell_j|\le2|t|$ and $|\ell_j^2-t^2|\le7|t|^3$. Then \eqref{eq:hdp-g1} gives
$|g_j-\frac14t^2|\le2|t|^3$, and with $|\ell|^2=3t^2+|w|^2$ it bounds $\psi$.
Hence on $E$
\begin{equation}\label{eq:hdp-pair}
 \bigl|g_i+g_j-\tfrac12t^2\bigr|\le4|t|^3\quad(i\ne j),
 \qquad\bigl|\psi-\tfrac34t^2\bigr|\le2t^4.
\end{equation}
So on $E$ each right side of \eqref{eq:hdp-constraints} is $\frac12t^2$ up to
$4|t|^3$. Now let
$\mathcal T(\delta)=\{M:M_1,M_2\le\delta,\ M_1+M_2\ge-\delta\}$. This is the
set where $M_1,M_2,M_3\le\delta$. It is a triangle with corners
$(\delta,\delta)$, $(\delta,-2\delta)$ and $(-2\delta,\delta)$, so it has area
$\frac92\delta^2$ and lies in $\{|M|_\infty\le2\delta\}$. Let $D_t$ be the
section of $D\cap U$ at $t$ and put $\delta_\pm=\frac12t^2(1\pm8|t|)\ge0$,
which are $\frac12t^2\pm4|t|^3$. If $M\in D_t$, then $(t,M)\in E$, and
\eqref{eq:hdp-constraints} and \eqref{eq:hdp-pair} give
$M_1,M_2,M_3\le\delta_+$. Conversely, if $M\in\mathcal T(\delta_-)$, then
$|M|_\infty\le2\delta_-\le\min(t^2,\varepsilon)$, so $(t,M)\in E\cap U$, and
\eqref{eq:hdp-pair} gives \eqref{eq:hdp-constraints}. So
$\mathcal T(\delta_-)\subseteq D_t\subseteq\mathcal T(\delta_+)$, and
\begin{equation}\label{eq:hdp-area}
 \tfrac98t^4(1-8|t|)^2\le|D_t|\le\tfrac98t^4(1+8|t|)^2.
\end{equation}

\emph{Step 4: integration.} Let $J=\int_{D\cap U}e^{-a\psi}$. We integrate
first over the section and then over $t$ (the change of coordinates has
Jacobian $1$). By \eqref{eq:hdp-pair} and \eqref{eq:hdp-area}, $J$ lies
between the integrals over $|t|\le\varepsilon$ of
$\frac98t^4(1\mp8|t|)^2e^{-a(\frac34t^2\pm2t^4)}$. We substitute
$t=y/\sqrt a$. Then these are $a^{-5/2}$ times integrals over
$|y|\le\varepsilon\sqrt a$ whose integrands tend to $\frac98y^4e^{-3y^2/4}$ as
$a\to\infty$ and are at most $\frac92y^4e^{-y^2/2}$, since $8|t|<1$ and
$2y^4/a=2t^2y^2\le2\varepsilon^2y^2$. By dominated convergence and
$\int_{\R}y^4e^{-3y^2/4}dy=\frac34\sqrt\pi(\frac43)^{5/2}$,
\[
 J=a^{-5/2}\bigl(\sqrt{3\pi}+o(1)\bigr).
\]
By Step 2 the same holds for the integral over $D$. It remains to find $Z$. As
$\Gamma$ is log-convex, $\Gamma(a+\frac12)^2\le a\Gamma(a)^2$ and
$\Gamma(a+1)^2\le(a+\frac12)\Gamma(a+\frac12)^2$, so
$Z=2\sqrt{\pi/a}\,(1+O(1/a))$. Then \eqref{eq:hdp-integral} gives
\[
 \Prob_\omega(D)=\frac{a^{3/2}}{8\pi^{3/2}}\cdot a^{-5/2}\sqrt{3\pi}\,(1+o(1))
 =\frac{\sqrt3}{8\pi a}(1+o(1))=\frac{3\sqrt3}{8\pi\omega}(1+o(1)).
 \qedhere
\]
\end{proof}
